% GENERATED FILE. Edit canonical lessons, then run npm run generate:books.
% canonical-source-hash: fnv1a64:71e804e3e9107984
% canonical-lessons: HI-C154-garam, HI-C154-gila, HI-C154-sukha, HI-C154-bhari, HI-C154-halka

\chapter{Hot, Wet, Dry, Heavy, Light}
\label{ch:hi-hot-wet-dry-heavy-light}
\hlchaptermodality{154}

\begin{chapteropening}
\textbf{By the end of this chapter:} \emph{I can say hot, wet, dry, heavy and light.}

Complete the last lesson of chapter 154: I can say hot, wet, dry, heavy and light.
\end{chapteropening}

\section[\texorpdfstring{\hi{गरम} (garam)}{गरम (garam)}]{\hi{गरम} (garam) — hot, warm}

\label{lesson:HI-C154-garam}



\begin{quote}
Before the new one: say the Hindi for an office, then the Hindi for a bank.

Walk one \textbf{\hi{दिन}}: \textbf{\hi{दोपहर}}, two \emph{pahars}, now the whole afternoon, then \textbf{\hi{रात}}, and \textbf{\hi{आधी} \hi{रात}}, half the night.
\end{quote}

\subsection*{You'll want to know: \hi{गरम}}
\textbf{\hi{गरम}} — \emph{garam} — \textquotedblleft{}hot, warm\textquotedblright{}.

\begin{cousinweb}
\textbf{\hi{गरम}} (\emph{garam}) is a loan, usually traced to Persian \emph{garm}, \textquotedblleft{}warm, hot\textquotedblright{}.
\end{cousinweb}

\begin{grammarlens}[title={What you've built}]
One more word for this chapter.
\end{grammarlens}

\subsection*{Your turn}
\begin{itemize}
\raggedright
  \item \emph{Say it:} \emph{garam}
  \item \emph{Say it:} \emph{garam}, once more
  \item \emph{Say it:} say \emph{garam}
  \item \emph{Recall:} say the Hindi for an office, then the Hindi for a bank, then say \emph{garam} again
\end{itemize}

\begin{tcolorbox}[breakable,colback=teal!4,colframe=teal!35!black,title={Before you move on}]
What does \hi{गरम} mean? (\textquotedblleft{}Hot, warm\textquotedblright{}.) Say it once more.
\end{tcolorbox}
\section[\texorpdfstring{\hi{गीला} (gīlā)}{गीला (gīlā)}]{\hi{गीला} (gīlā) — wet}

\label{lesson:HI-C154-gila}



\begin{quote}
Before the new one: say the Hindi for a bank, then the Hindi for hot.
\end{quote}

\subsection*{You'll want to know: \hi{गीला}}
\textbf{\hi{गीला}} — \emph{gīlā} — \textquotedblleft{}wet\textquotedblright{}.

\begin{cousinweb}
\textbf{\hi{गीला}} (\emph{gīlā}) is a tadbhava, a word inherited and reshaped by sound change. It comes from \emph{*grilla}, \textquotedblleft{}wet, damp\textquotedblright{}, a form CDIAL reconstructs from its descendants (CDIAL 4386). CDIAL calls a further link with \emph{*gṛdla} very doubtful.
\end{cousinweb}

\begin{grammarlens}[title={What you've built}]
One more word for this chapter.
\end{grammarlens}

\subsection*{Your turn}
\begin{itemize}
\raggedright
  \item \emph{Say it:} \emph{gīlā}
  \item \emph{Say it:} \emph{gīlā}, once more
  \item \emph{Say it:} say \emph{gīlā}
  \item \emph{Recall:} say the Hindi for a bank, then the Hindi for hot, then say \emph{gīlā} again
\end{itemize}

\begin{tcolorbox}[breakable,colback=teal!4,colframe=teal!35!black,title={Before you move on}]
What does \hi{गीला} mean? (\textquotedblleft{}Wet\textquotedblright{}.) Say it once more.
\end{tcolorbox}
\section[\texorpdfstring{\hi{सूखा} (sūkhā)}{सूखा (sūkhā)}]{\hi{सूखा} (sūkhā) — dry}

\label{lesson:HI-C154-sukha}



\begin{quote}
Before the new one: say the Hindi for hot, then the Hindi for wet.
\end{quote}

\subsection*{You'll want to know: \hi{सूखा}}
\textbf{\hi{सूखा}} — \emph{sūkhā} — \textquotedblleft{}dry\textquotedblright{}.

\begin{cousinweb}
\textbf{\hi{सूखा}} (\emph{sūkhā}) is a tadbhava, a word inherited from Sanskrit and reshaped by sound change. It comes from \emph{śuṣka}, \textquotedblleft{}dried\textquotedblright{} (CDIAL 12548).
\end{cousinweb}

\begin{grammarlens}[title={What you've built}]
One more word for this chapter.
\end{grammarlens}

\subsection*{Your turn}
\begin{itemize}
\raggedright
  \item \emph{Say it:} \emph{sūkhā}
  \item \emph{Say it:} \emph{sūkhā}, once more
  \item \emph{Say it:} say \emph{sūkhā}
  \item \emph{Recall:} say the Hindi for hot, then the Hindi for wet, then say \emph{sūkhā} again
\end{itemize}

\begin{tcolorbox}[breakable,colback=teal!4,colframe=teal!35!black,title={Before you move on}]
What does \hi{सूखा} mean? (\textquotedblleft{}Dry\textquotedblright{}.) Say it once more.
\end{tcolorbox}
\section[\texorpdfstring{\hi{भारी} (bhārī)}{भारी (bhārī)}]{\hi{भारी} (bhārī) — heavy}

\label{lesson:HI-C154-bhari}



\begin{quote}
Before the new one: say the Hindi for wet, then the Hindi for dry.
\end{quote}

\subsection*{You'll want to know: \hi{भारी}}
\textbf{\hi{भारी}} — \emph{bhārī} — \textquotedblleft{}heavy\textquotedblright{}.

\begin{cousinweb}
\textbf{\hi{भारी}} (\emph{bhārī}) is a tadbhava, a word inherited from Sanskrit and reshaped by sound change. It comes from \emph{bhārika}, \textquotedblleft{}forming a load, heavy\textquotedblright{} (CDIAL 9465). It is built on \emph{bhāra}, \textquotedblleft{}load\textquotedblright{}.
\end{cousinweb}

\begin{grammarlens}[title={What you've built}]
One more word for this chapter.
\end{grammarlens}

\subsection*{Your turn}
\begin{itemize}
\raggedright
  \item \emph{Say it:} \emph{bhārī}
  \item \emph{Say it:} \emph{bhārī}, once more
  \item \emph{Say it:} say \emph{bhārī}
  \item \emph{Recall:} say the Hindi for wet, then the Hindi for dry, then say \emph{bhārī} again
\end{itemize}

\begin{tcolorbox}[breakable,colback=teal!4,colframe=teal!35!black,title={Before you move on}]
What does \hi{भारी} mean? (\textquotedblleft{}Heavy\textquotedblright{}.) Say it once more.
\end{tcolorbox}
\section[\texorpdfstring{\hi{हल्का} (halkā)}{हल्का (halkā)}]{\hi{हल्का} (halkā) — light (in weight)}

\label{lesson:HI-C154-halka}



\begin{quote}
Before the new one: say the Hindi for dry, then the Hindi for heavy.
\end{quote}

\subsection*{You'll want to know: \hi{हल्का}}
\textbf{\hi{हल्का}} — \emph{halkā} — \textquotedblleft{}light (in weight)\textquotedblright{}.

\begin{cousinweb}
\textbf{\hi{हल्का}} (\emph{halkā}) is a tadbhava, a word inherited from Sanskrit and reshaped by sound change. It comes from \emph{laghu}, \textquotedblleft{}light\textquotedblright{} (CDIAL 10896). CDIAL notes that its \emph{l} and \emph{h} swapped places and that it took an added \emph{-k-}.
\end{cousinweb}

\begin{grammarlens}[title={What you've built}]
One more word for this chapter.
\end{grammarlens}

\subsection*{Your turn}
\begin{itemize}
\raggedright
  \item \emph{Say it:} \emph{halkā}
  \item \emph{Say it:} \emph{halkā}, once more
  \item \emph{Say it:} say \emph{halkā}
  \item \emph{Recall:} say the Hindi for dry, then the Hindi for heavy, then say \emph{halkā} again
\end{itemize}

\begin{tcolorbox}[breakable,colback=teal!4,colframe=teal!35!black,title={Before you move on}]
What does \hi{हल्का} mean? (\textquotedblleft{}Light (in weight)\textquotedblright{}.) Say it once more.
\end{tcolorbox}
